\section{Temperaturverhalten am Prüfstand}
\label{sec:temperatur}

Vor den Laufversuchen wurde am Rad 97351 die auf der Zahnflanke gemessene Zahnmassentemperatur über dem Drehmoment und der Drehzahl nach \secref{sec:auswertung_laufversuche} aufgenommen. Ziel war die Einstellung der Öltemperatur für die Zahnmassentemperatur von $80$\,\textdegree C. Die Öltemperatur an der Messstelle vor der Prüfverzahnung betrug $75$\,\textdegree C. \figref{fig:temperatur_matrix_97351} zeigt die sechs Werte als Balken gegen den Bezug von $80$\,\textdegree C.

% TK erledigt (nur noch ME-Nummern): Die Bezeichnungen der Excel mit kst-E_ bitte überall rausnehmen und immer nur mit der ME Nummer arbeiten.

\begin{figure}[htbp]
    \centering
    \includegraphics[width=\textwidth]{kap4/temperatur_matrix_97351}
    \caption{Zahnflankentemperatur nach zwei Stunden am Rad 97351 über dem Raddrehmoment bei $2250$ und $3000$\,1/min am Rad, Öltemperatur $75$\,\textdegree C an der Messstelle, Wärmebildkamera mit Emissionsgrad $1$, Bezug $80$\,\textdegree C.}
    \label{fig:temperatur_matrix_97351}
\end{figure}

% TK (folgt mit Paket B3, Bild wird neu gezeichnet): Den Balken fehlt ein Muster für schwarz weiß und auch hier weider das Grid in den Hintergrund bringen. Balken gerne auch etwas dünner. Sieht so platzeinnehmend aus, obwohl die Grafik am Ende wenig aussagt, also nicht minimalprinzip meiner Meinung nach.

\figref{fig:waermebild_16nm_3000} zeigt als Beispiel die Aufnahme bei $16$\,Nm und $3000$\,1/min mit dem Messpunkt auf der Zahnflanke. Die übrigen fünf Aufnahmen stehen in \secref{sec:anhang_h}.

\begin{figure}[htbp]
    \centering
    \includegraphics[width=0.8\textwidth]{kap4/waermebild_16nm_3000}
    \caption{Wärmebildaufnahme des Rades 97351 nach zwei Stunden bei $16$\,Nm und $3000$\,1/min am Rad, Emissionsgrad $1$. Der Messpunkt liegt auf der Zahnflanke, die Legende nennt die Temperaturspanne der Aufnahme.}
    \label{fig:waermebild_16nm_3000}
\end{figure}

Bei $8$ und $12$\,Nm liegt die Flankentemperatur mit $79{,}5$ bis $81{,}9$\,\textdegree C innerhalb von $2$\,K um den Bezug. Bei $16$\,Nm steigt sie auf $84{,}1$\,\textdegree C bei $2250$\,1/min und auf $87{,}8$\,\textdegree C bei $3000$\,1/min. Die höhere Drehzahl trägt zwischen $0{,}9$\,K bei $12$\,Nm und $3{,}7$\,K bei $16$\,Nm bei. Der Schritt von $8$ auf $16$\,Nm bringt bei $3000$\,1/min $6{,}3$\,K. Die Öltemperatur von $75$\,\textdegree C hält den Bezug damit auf den Stufen $8$ und $12$\,Nm, auf der Stufe $16$\,Nm läuft der Zahn wärmer. In den Läufen bei $16$\,Nm wurde die Solltemperatur an der Messstelle deshalb auf $67$\,\textdegree C gesetzt (vgl. \tabref{tab:versuchsprogramm_pa46}). Maßgebend ist die Zahnmassentemperatur von $80$\,\textdegree C, auf die alle Versuche an dieser Verzahnung bezogen sind, auch die anderer Lehrstühle. Eine Kastentemperatur wurde in der Matrix nicht aufgezeichnet. Nach der Matrix lief das Rad bei $12$\,Nm weiter und versagte nach $2{,}42 \cdot 10^{6}$ Lastwechseln über alle Stufen mit den vollständig abgebrochenen Zähnen 34 bis 41 (\secref{sec:schadensbilder}). Dieser Lauf trägt eine gestufte Vorgeschichte und erscheint in \secref{sec:woehlerlinien} nur als Marker.

% Herkunft:
% Absatz 1 und 2 (Matrix): 70_input/72_Kleingetriebepruefstand_KV3/kst-E_Laufversuche_KV3.xlsx, Zeilen 77 bis 83
%   (kst-E_63, Rad 97351, Ritzel 95919, Öltemperatur 75, Ziel 80, Ritzelmoment 11,8/7,9/7,8/15,7 Nm, Drehzahl
%   2294 und 3058 1/min am Ritzel = 2250 und 3000 am Rad, Bemerkung "nach 2h ... °C Zahnflankentemperatur mit
%   Emissionsgrad 1"); 80_output/messungen/woehler/vergleich.md Tabelle "Temperaturmatrix kst-E_63" (2250:
%   79,5/81,0/84,1; 3000: 81,5/81,9/87,8 °C); Bild temperatur_matrix_97351 (measurements.py woehler, L4).
%   Wärmebildkamera, Öltemperatur = Messstelle vor der Prüfverzahnung, Weiterlauf bei 12 Nm bis zum Ende:
%   Nutzer 2026-10-18 (ADR-118 Amendment vom Abend, Punkt 4 und 5). Differenzen: Drehzahl 2,0/0,9/3,7 K,
%   Drehmoment bei 3000 1/min 6,3 K. 67 °C der Läufe 70 bis 73: Spalte G der Mappe. Endzahl 2 424 628:
%   Spalte S Zeile 83; "5-6 kaputte Zähne": Spalte T Zeile 83. VDI 2736 Blatt 4 Abschnitt 8.2: über 3.8.
%   67 °C und Emissionsgrad: Nutzer 2026-10-18 (Vorversuche zeigten bei erhöhtem Moment deutlich höhere
%   Zahnmassentemperaturen; Vereinbarung mit anderen Lehrstühlen auf 80 °C; Emissionsgrad 1 aus einem Vorversuch
%   der Forschungsstelle mit Messstreifen); Wärmebilder: 70_input/72_Kleingetriebepruefstand_KV3/Bilder/
%   Waermebildaufnahme_Kunststoffraeder/<T>_<n>.png, geschnitten mit scripts/thesis_thermal_images.py.
% Bildwunsch: kap4/temperatur_matrix_97351, kap4/waermebild_16nm_3000 (bilder.yaml); die übrigen fünf in Anhang H.
% Bibwunsch: keiner.
% Nomenklatur: keine neuen Zeichen.
% Offen: Kastentemperatur der Matrix (nicht aufgezeichnet).
